%---------------------%
% Temporary registers %
%---------------------%

\newbox   \tmpboxA
\newbox   \tmpboxB
\newcount \tmpcountA
\newcount \tmpcountB
\newdimen \tmpdimenA
\newdimen \tmpdimenB

%--------------------------%
% Print/electronic variant %
%--------------------------%

\newif \ifprintversion
\printversiontrue

%------------------%
% Basic shorthands %
%------------------%

% Shorthand for \expandafter
\let \ea=\expandafter
% Un-\def a control sequence or active character
\def \undef #1{\let#1\undefined}
% Define control sequence with name specified in <balanced text>
\def \defbyname #1{\ea\def\csname#1\endcsname}

%------------------%
% \foreach, \xargs %
%------------------%

% Require `{` after \foreach
\def \foreach #{\foreachA}

\def \foreachA #1#2{\foreachB#2#1\foreachC\foreachD}
\def \foreachB #1#2#3\foreachD {\ifx #2\foreachC \else #1{#2}\foreachB #1#3\foreachD \fi}
% \foreachC and \foreachD are "sentinels" that should never be expanded:
\def \foreachC {\error{\string\foreachC has been expanded!}}
\def \foreachD {\error{\string\foreachD has been expanded!}}

\def \xargs #1#2;{\foreachB#1#2\foreachC\foreachD}

%---------------------%
% \scantoeol, \eoldef %
%---------------------%

\def \scantoeol {\bgroup \catcode`\^^M=13 \scantoeolA}
{\catcode`\^^M=13 \gdef\scantoeolA#1#2^^M{\egroup #1{#2}}}
\def \eoldef #1#2#{%
  \edef#1{\noexpand\scantoeol \ea\noexpand\csname _eoldef:\csstring#1\endcsname}
  \ea\def \csname _eoldef:\csstring#1\endcsname#2}

%------------------%
% \ignorespacetoks %
%------------------%

%
\def \ignorespacetoks #1{\let\afterignore#1\ignorespacetoksA}
%
\def \ignorespacetoksA {\futurelet\tmptokenA\ignorespacetoksB}
% \ignorespacetoksB decides what to do based on what \tmptokenA is
\def \ignorespacetoksB {%
  \ifcat\space\tmptokenA \ea\ignorespacetoksC
  \else \ea\afterignore \fi}
% \ignorespacetoksC gobbles the space and "inserts" \ignorespacetoksA token
\def \ignorespacetoksC {\afterassignment\ignorespacetoks \let\tmptokenA= }

%-------------------------%
% \readopt, \readoptional %
%-------------------------%

% \readopt loads next token and delegates main work to \readoptA
\def \readopt #1{\let\afteropt#1\futurelet\tmptokenA\readoptA}
% \readoptA decides what to do based on what \tmptokenA is
\def \readoptA {\ifx[\tmptokenA \ea\readoptB \else \ea\readoptC \fi}
% \readoptB loads the optional parameter
\def \readoptB [#1]{\def\theoptional{#1}\afteropt}
% \readoptC undefines \theoptional and expands to \afteropt
\def \readoptC {\undef\theoptional \ignorespacetoks\afteropt}

\let \readoptional=\readopt

%----------------%
% \protecttokens %
%----------------%

% The \protecttokens macro converts given token sequence to sequence suitable
% to be written to file using \write by doing following:
%   (1) \catcode of begin-group and end-group character and space tokens (left
%       and right curly brace, ASCII space, ASCII horizontal tabulator) changes
%       to 12 (other).
%   (2) \string primitive is applied to each token of token list trasformed as
%       specified by (1).
%   (3) The final token list will be saved into \protectedtokens.

\def \protecttokens #1{
  \bgroup
  \catcode`\{=12 \catcode`\}=12
  \catcode`\ =12 \catcode`\^^I=12
  \xdef\protectedtokens {\ea\foreach\ea{\scantextokens{#1}}\protecttoken}%
  \egroup}
\def \protecttoken #1{\string#1}

%---------------------------%
% Flavors of \the primitive %
%---------------------------%

% \Athe and \athe take \count register as its argument and converts it into
% a letter: 1-->"A"/"a", 2-->"B"/"b", ..., 26-->"Z"/"z", "??" otherwise.
\def \athe #1{\ifcase#1??\or
  a\or b\or c\or d\or e\or f\or g\or h\or i\or j\or k\or l\or m\or
  n\or o\or p\or q\or r\or s\or t\or u\or v\or w\or x\or y\or z\else ??\fi}
\def \Athe #1{\ifcase#1??\or
  A\or B\or C\or D\or E\or F\or G\or H\or I\or J\or K\or L\or M\or
  N\or O\or P\or Q\or R\or S\or T\or U\or V\or W\or X\or Y\or Z\else ??\fi}

%----------%
% Messages %
%----------%

% Warnings
\def \warn #1{\message{^^J>>> WARNING! #1^^J}}
% Debug messages
\def \debugmsg #1{\message{^^J>>> DEBUG: #1^^J}}
% Messages intended for tracing user errors
\newcount\showtracemessages
\def \tracemsg #1{\ifnum\showtracemessages>\z@ \message{^^J>>> TRACE: #1^^J}\fi}
\showtracemessages=1

%-------------------------------------%
% Redefinition of glue macros, \?, \. %
%-------------------------------------%

\newdimen \thinspaceamount       \thinspaceamount=0.16667em
\newdimen \verythinspaceamount   \verythinspaceamount=0.06667em

\def \negthinspace      {\ifmmode \mskip-\thinmuskip \else \kern-\thinspaceamount \fi}
\def \thinspace         {\ifmmode \mskip \thinmuskip \else \kern \thinspaceamount \fi}
\def \negverythinspace  {\kern-\verythinspaceamount}
\def \verythinspace     {\kern \verythinspaceamount}

\let \!=\negthinspace
\let \,=\thinspace
\let \?=\negverythinspace
\let \.=\verythinspace

%-------%
% Logos %
%-------%

\def \TeX {T\kern -0.1667em\relax \lower 0.5ex\hbox{E}\kern -0.125em\relax X}
\def \pdfTeX {pdf\TeX}
\def \LaTeX {L\kern 0.16em\llap{\raise 0.20em\hbox{\the\scriptfont0 A}}\kern -0.16em\TeX}
\def \OpTeX {Op\kern -0.1em\relax \TeX}
\def \LuaTeX {Lua\TeX}
\def \XeTeX {X\kern -0.125em\relax
  \phantom E\pdfsave \rlap{\pdfsetmatrix{-1 0 0 1}\lower0.5ex\hbox{E}}\pdfrestore
  \kern -0.1667em\relax \TeX}

%------------------------------------------------%
% \include, \includedir, \includesub, \softinput %
%------------------------------------------------%

\let \currentdir=\empty
% Used by \softinput to test file existence
\newread \softinputread
\newif \iffilemissing

\def \includesub #1:#2 {{%
  \edef\includesubZ{#1}%
  \ifx \includesubZ\empty \else
    \edef\currentdir{\currentdir\includesubZ/}%
  \fi \undef\includesubZ \input{\currentdir#2}}}

\def \includedir #1 {\includesub#1:main }

\def \include #1 {\includesub:#1 }

\def \softinput#1 {%
  \openin \softinputread=#1
  \ifeof \softinputread
    \def\softinputZ{\undef\softinputZ \filemissingtrue}%
  \else
    \closein\softinputread
    \def\softinputZ{\undef\softinputZ \filemissingfalse \input{#1}}%
  \fi \softinputZ}

%----------%
% \mathbox %
%----------%

\def \mathbox #{\mathboxA}
\def \mathboxA #1{\mathchoice{\mathboxT{#1}}{\mathboxT{#1}}{\mathboxS {#1}}{\mathboxSS{#1}}}
\def \mathboxT #1{\hbox{#1}}
\def \mathboxS #1{\hbox{\currfontsize=0.7\currfontsize \selectfont\selectmathfont #1}}
\def \mathboxSS #1{\hbox{\currfontsize=0.5\currfontsize \selectfont\selectmathfont #1}}

%-----------------------------------%
% Other utilitity macros for layout %
%-----------------------------------%

\def \skiplines #1 {%
  \tmpdimenA=#1\baselineskip \relax
  \ifdim\lastskip<\tmpdimenA \removelastskip \vskip\tmpdimenA \fi}

\long\def \ignorepar {\futurelet\ignoreparZ\ignoreparA}
\long\def \ignoreparA {%
  \ifx \ignoreparZ\par \ea\ignoreparB
  \else \undef\ignoreparZ \fi}
\long\def \ignoreparB #1{\ignorepar}

% Taken from OpTeX source code
\def \firstnoindent {%
  \global\everypar{%
    \global\everypar={}%
    \setbox7=\lastbox}}

% Redefinition of plain TeX's \strut, new scales with baselineskip
\def \strut {\vrule
    height 0.75\baselineskip
    depth 0.25\baselineskip
    width \z@ \relax}
